\documentclass[11pt]{article}
\usepackage[margin=1in]{geometry}
\usepackage{booktabs}
\title{Novelty Detection and Threshold Calibration}
\author{Devis W. Saputra}
\begin{document}
\maketitle

\begin{abstract}
Digit 0 is treated as an unseen class while an Isolation Forest is fitted only on digits 1 through 9. A separate normal-only validation set calibrates thresholds at several false-positive budgets before evaluation on a mixed held-out test set.
\end{abstract}

\section{Method}
The normal data are split into 841 model-fitting and 211 threshold-validation examples. The mixed test set contains 629 examples. No digit-0 case is used for fitting or threshold calibration.

\section{Ranking}
The continuous anomaly score reaches ROC-AUC 0.8071 and Average Precision 0.2458.

\section{Operating Points}
\begin{tabular}{lrrrrr}
\toprule
Validation FPR & Test FPR & Precision & Recall & F1 & Bal. Acc. \\
\midrule
0.05 & 0.0511 & 0.2750 & 0.1774 & 0.2157 & 0.5631 \\
0.10 & 0.1129 & 0.2644 & 0.3710 & 0.3087 & 0.6290 \\
0.15 & 0.1834 & 0.2571 & 0.5806 & 0.3564 & 0.6986 \\
\bottomrule
\end{tabular}

\section{Interpretation}
Increasing the tolerated false-positive budget improves novelty recall. Ranking quality and threshold choice are separate engineering decisions, and the operating point should reflect downstream review costs.

\section{Limitations}
Digit 0 is a benchmark novelty, not a realistic open-world anomaly distribution. Future work should test multiple unseen classes, drift, and external out-of-distribution data.

\end{document}
